\dish{Red Pepper and Doubanjiang Pasta}
\label{pasta-poivron-doubanjiang}
\altdish{Pepper pasta with doubanjiang}
\altdish{豆瓣酱 and red pepper pasta sauce}
\serves{2}
\prep{35 min}
\sourcep{chatgpt}

\begin{ingredients}
  \ingr{}{}{butter}
  \ingr{2}{}{leek (white and light green), thinly sliced, or 2 onions}
  \ingr{3--6}{}{red bell peppers, chopped}
  \ingrS{1}{clove}{garlic, minced}
  \ingr{1--2}{tsp}{doubanjiang}
  \ingrS{160--180}{g}{pasta}
  \ingrS{}{}{pasta water}
  \ingr{}{}{lemon juice or wine vinegar, to taste}
  \ingr{}{}{salt}
\end{ingredients}


\begin{recipe}
  \begin{enumerate}

  \item Melt 2~Tbsp butter over medium heat.  Sweat the leek or onion
    until soft, 5--7~min.

  \item Add the peppers.  Cook, stirring occasionally, until soft and
    starting to concentrate, 15--20~min.

  \item Push vegetables aside.  Add doubanjiang to the exposed
    butter and fry 30--60~sec until fragrant.  Stir in garlic for
    the last 15~sec.  Mix everything together.

  \item Cook pasta in salted water.  Reserve a mugful of pasta
    water before draining.

  \item Blend the pepper mixture smooth with a splash of pasta
    water and the remaining 1~Tbsp butter.

  \item Toss the sauce with the pasta, loosening with pasta water
    to coat.  Sharpen with a little lemon juice or vinegar if the
    peppers taste too sweet.  Taste before salting --- the
    doubanjiang carries most of it.

  \end{enumerate}
\end{recipe}

\notes{} Serve with \hyperref[carrot-tomato-vinaigrette]{Carrot and
  Tomato Salad}, kept sharply acidic to contrast this dish's sweetness
and richness.
